% ED26-REQ-N-c-004-01: original Fagin 1974 presentation.
% Source identity and authoring departures are recorded outside the reader text.

\subsection{Fagin 1974: generalized spectra and recognizability}
\label{sec:p09-n-c-004-01}
\index{Fagin's theorem}\index{generalized spectrum}

This presentation fixes the original finite-structure encoding before stating
Theorem~6 of \cite{Fagin1974}.  The public modern capture statement remains
Theorem~\ref{thm:p03-fagin}; the two clauses below retain the original paper's
different resource measures.  In particular, a finite \emph{empty vocabulary}
does not mean an empty universe.

\begin{hbdefinition}{def:p09-fagin74-spectrum}{Expansion, reduct and generalized spectrum}
A finite similarity type is a finite list of relation and function symbols
with specified arities.  Relation arities are positive; constants are
zero-ary functions.  Fagin's finite structures have nonempty finite universes.
Write \(\operatorname{Fin}(S)\) for their class at type \(S\).
For disjoint finite types \(S,T\), an \(S\cup T\)-structure
\(\mathcal B\) expands \(\mathcal A\) when it has the same universe and
the same interpretations of every symbol in \(S\); its \(S\)-reduct is
\(\mathcal B\mathbin{\upharpoonright}S=\mathcal A\).
For a first-order \(S\cup T\)-sentence \(\sigma\), put
\[
 \operatorname{Mod}_{S}\exists T\sigma
 =\{\mathcal A\in\operatorname{Fin}(S):
       \text{some expansion }\mathcal B\text{ satisfies }\sigma\}.
\]
Here \(\exists T\) existentially quantifies all the relation and function
symbols of \(T\); only the symbols of \(S\) remain free at second order.
Such a class is an \((S,T)\)-spectrum.  It is an \(S\)-spectrum if some
finite \(T\) works, and a generalized spectrum if some \(S\) works.
When \(S=\varnothing\), identify the class of pure finite sets with its set
of positive cardinalities; this is a numerical spectrum.  In a monadic
generalized spectrum the auxiliary type \(T\) consists only of unary
relations \cite[\S2, pp.~47--49]{Fagin1974}.
\end{hbdefinition}

\begin{hbdefinition}{def:p09-fagin74-encoding}{The original encoding and the two input lengths}
For Theorem~6 the input type \(S=\{P_1,\ldots,P_r\}\) is fixed and
contains only relations, of positive arities \(m_1,\ldots,m_r\), in a fixed
order.  On universe \(\{1,\ldots,n\}\), let \(b_i\) list the truth values
of \(P_i\) on all \(n^{m_i}\) argument tuples in lexicographic order.
With \(a\) the usual binary representation of \(n\), the encoding is
\[
 e(\mathcal A)=a\#b_1\#b_2\#\cdots\#b_r,
 \qquad
 E(\mathcal C)=\{e(\mathcal A):\mathcal A\in\mathcal C,
                             \dom(\mathcal A)=\{1,\ldots,n\}\}.
\]
For \(r=0\) the numerical input is the binary representation of \(n\).
A constant separator convention does not change the time classes.
Put \(\ell=\operatorname{len}(e(\mathcal A))\).  If \(r>0\), then
\[
 \ell=\lfloor\log_2n\rfloor+1+r+\sum_{i=1}^{r}n^{m_i},
 \qquad n\leq\ell\leq C_S n^{\max_i m_i}\quad(n\geq2)
\]
for a fixed constant \(C_S\).  If \(r=0\), then
\(\ell=\lfloor\log_2n\rfloor+1\), so
\(2^{\ell-1}\leq n<2^\ell\).
The source's \(NP\) uses time \(\ell^k\) for some fixed positive integer
\(k\).  Its \(NP_1\) uses time \(2^{k\ell}\) for some such \(k\),
equivalently time polynomial in the \emph{value} \(n\) for binary numerical
inputs.  This \(NP_1\) is the source's exponential-time notation, not its
notation for polynomial time on a unary input
\cite[\S3, p.~51; \S4, pp.~52--53]{Fagin1974}.
\end{hbdefinition}

\begin{hbtheorem}{thm:p09-fagin74-theorem6}{proved=annex:p9.n-c-004.1.theorem6}{Fagin's original Theorem 6}
Fix the relational type \(S\) and encoding of
Definition~\ref{def:p09-fagin74-encoding}.  Let
\(\mathcal C\subseteq\operatorname{Fin}(S)\) be closed under isomorphism.
\begin{enumerate}
 \item If \(S\ne\varnothing\), then \(\mathcal C\) is an
       \(S\)-spectrum if and only if \(E(\mathcal C)\in NP\).
 \item If \(S=\varnothing\), then its set of cardinalities is a spectrum
       if and only if \(E(\mathcal C)\in NP_1\).
\end{enumerate}
In either case, an equivalent unified formulation is the existence of a
nondeterministic multitape machine recognizing exactly \(E(\mathcal C)\)
which accepts each member along some computation of at most \(n^k\) steps
for some fixed positive integer \(k\), where \(n\) is the represented universe size
\cite[Theorem~6 and its note, p.~53]{Fagin1974}.
\end{hbtheorem}

\paragraph{The two translations.}
For a spectrum, first represent auxiliary functions by total single-valued
relation graphs and put the first-order matrix in prenex form.  The machine
rejects malformed encodings, guesses the auxiliary characteristic tables once,
and evaluates the fixed matrix, respecting each first-order universal or
existential quantifier.  Enumerating all assignments is polynomial in \(n\)
because the sentence, arities and number of quantified variables are fixed.
For a nonempty input type this yields a polynomial-in-\(\ell\) bound;
in the empty-type case the same argument yields the \(NP_1\) bound.

For the converse, Fagin supplies a construction rather than merely invoking
a machine-class slogan.  His modified machine receives \(r+1\) tapes:
the reversed binary representation of \(n\) and the separate relation words
\(b_i\).  An existentially chosen order numbers the original universe;
\(k\)-tuples number time instants and tape squares.  The auxiliary symbols
are the order \(<\), its endpoints \(c_1,c_2\), its cyclic successor
function \(S\), a tuple-successor relation \(S_1\), the state function
\(q\), tape-symbol functions \(v_i\), head relations \(H_i\), and a
binary-digit function \(G\).  The source's nine conjunct groups specify:
symbol/state codes; the order; tuple successor; binary digits; initial/final
state and unique heads; the numerical input tape; the relation input tapes;
unchanged final configurations; and a disjunction of legal transitions for
each table entry \cite[pp.~54--58]{Fagin1974}.

The full reconstruction in Annex~\ref{annex:p9.n-c-004.1.theorem6} includes
the meaning and first-order definition of all these blocks, both soundness and
completeness of the tableau, the small-universe finite modification, and the
function-to-relation conversion.  It uses explicit tape/state range constraints
and the correct carry condition for binary increment; these are part of the
reconstruction, not a claim that the source prints every necessary condition
in that form.  Fagin describes the guessing machine informally and leaves
prenex conversion, graph coding and finite modification as standard steps;
the annex supplies those dependencies at the precision needed here.

\begin{hbexample}{def:p09-fagin74-even-witness}{Even size as an existential certificate}
The source's introductory spectrum example is
\[
 \exists f\,\bigl(\forall x\,f(x)\ne x\ \land
                  \forall x\,f(f(x))=x\bigr).
\]
Its witness is a fixed-point-free involution, hence a partition into pairs.
This defines precisely the positive even cardinalities
\cite[pp.~43--44]{Fagin1974}.  For a purely relational ESO sentence use
\(F(x,y)\) with
\[
 \forall x\exists!y\,F(x,y),\qquad
 \forall x\,\neg F(x,x),\qquad
 \forall x\forall y\,(F(x,y)\ \longleftrightarrow\ F(y,x)).
\]
Existentially quantify \(F\) and conjoin these three first-order conditions.
All three conditions are needed: irreflexivity alone does not partition the
universe, and an arbitrary permutation without fixed points can have a
three-cycle.  The example is unchanged by adding unused input relations.
\end{hbexample}

\begin{hblemma}{thm:p09-fagin74-unordered-boundary}{proved=annex:p9.n-c-004.1.unordered-boundary}{The unordered parity boundary}
On nonempty finite pure sets, even cardinality is decidable in polynomial time
under an explicit universe representation and is definable in ESO, but is not
definable by an order-free first-order least-fixed-point sentence.
The last assertion also holds on structures in any fixed relational type when
all its input relations are empty.
\end{hblemma}

This is an editorial comparison, not a fixed-point theorem claimed for the
1974 article.  In pure sets, a fixed formula and all its fixed-point stages
can distinguish only finitely many equality patterns of tuples; beyond a
formula-dependent size they behave identically on all universes.
The annex proves this invariance, including parameters and the iteration
stages.  Even size alternates forever and therefore fails that test.
The order in Theorem~\ref{thm:p04-immerman-vardi} is consequently not
removable.  In Fagin's theorem the order is an existentially guessed witness,
not a distinguished relation which the input must provide.

\begin{table}[htbp]
\centering\small
\begin{tabular}{@{}p{.23\textwidth}p{.32\textwidth}p{.35\textwidth}@{}}
\toprule
Reasoning task & What varies and how it is represented & Correct conclusion \\
\midrule
Recognize a fixed ESO property, \(S\ne\varnothing\) & Input structure varies;
fixed finite relational type; original characteristic-table encoding &
Exactly the isomorphism-closed \(NP\) properties, by original Theorem~6(1). \\
Recognize a numerical spectrum & Binary universe size varies; input type empty
& Original Theorem~6(2) gives \(NP_1\), not \(NP\) in the binary length. \\
Evaluate a fixed first-order matrix after guessing its fixed auxiliary relations
& Structure and guessed tables vary; formula and arities fixed & Deterministic
checking is polynomial in \(n\); one guess serves all assignments. \\
Ask whether ordered fixed-point logic captures deterministic polynomial time
& Structure varies and its total order is part of the input vocabulary &
Use Theorem~\ref{thm:p04-immerman-vardi}; dropping the order fails on parity. \\
Evaluate an input formula & Formula varies, either with the structure fixed or
with both formula and structure varying & Expression or combined complexity is
a different question; Fagin's fixed-property theorem supplies no uniform
formula-input polynomial bound. \\
\bottomrule
\end{tabular}
\caption{Task and input axes for the 1974 result and its editorial boundaries.}
\label{tab:p09-fagin74-input-axes}
\end{table}

\begin{hbwarning}{def:p09-fagin74-encoding-boundary}{Representation, invariance and order are separate conditions}
An arbitrary language of strings need not be invariant under relabeling the
represented structure.  Conversely, two isomorphic structures can have
different characteristic words.  Isomorphism closure is what makes the
guessed ordering harmless in the tableau proof.  Changing to a succinct
representation of a relation can change \(\ell\) exponentially and is not
authorized by Theorem~6.  Changing the empty-type size input from binary
\(n\) to the unary word \(1^n\) does yield an \(NP\) reformulation,
because both the encoding conversion and the \(n^k\) simulation then cost
polynomial time in the new input length.  That is a derived reformulation,
not the encoding printed in the original theorem.
\end{hbwarning}
